% Generated draft: requires mathematical and source review.
\documentclass[11pt]{article}
\usepackage[margin=25mm]{geometry}
\usepackage{fontspec}
\setmainfont{Latin Modern Roman}
\usepackage{amsmath,amssymb,mathtools,unicode-math}
\setmathfont{Latin Modern Math}
\usepackage{enumitem,needspace,xurl,hyperref}
\setlength{\parindent}{0pt}
\setlength{\parskip}{5pt}
\setlength{\emergencystretch}{3em}
\newcommand{\N}{\mathbb{N}}
\newcommand{\Z}{\mathbb{Z}}
\newcommand{\Q}{\mathbb{Q}}
\newcommand{\R}{\mathbb{R}}
\newcommand{\C}{\mathbb{C}}
\newcommand{\F}{\mathbb{F}}
\newcommand{\A}{\mathbb{A}}
\newcommand{\PP}{\mathbb{P}}
\newcommand{\E}{\mathbb{E}}
\newcommand{\eps}{\varepsilon}
\newcommand{\dd}{\,\mathrm{d}}
\begin{document}
\section*{UnsolvedMath ID: 30000731}
\textbf{Source problem code:} catalogue entry
\subsection{Definitions and mathematical statement}
Let \(\mathcal H=(V,\mathcal E)\) be a finite hypergraph, with board \(V\) and winning sets \(\mathcal E\subseteq2^V\). In the usual weak-win or Maker--Breaker interpretation, Maker seeks to claim every vertex of some \(A\in\mathcal E\), while Breaker seeks to prevent that outcome. This is distinct from a strong game where both players have the same complete-set winning objective. A biased \((a,b)\) version typically lets Maker and Breaker claim respectively \(a\) and \(b\) previously unclaimed vertices per turn, subject to the convention for a final incomplete turn. The source does not specify the biases, starting player or full rules.

The first request is to extend Beck's weak-win Meta-Conjecture to biased games. The Meta-Conjecture itself is not stated, so its quantitative hypothesis, proposed strategy guarantee and bias dependence cannot be reconstructed from the name alone. No particular random-play threshold, potential sum or edge-size/degree inequality is substituted for the missing conjecture.

The second request is for useful positive results on sparse hypergraphs. The cleaned statement explicitly allows replacing complete occupation by a suitable majority or surplus objective, whereas the original asks more generally what remains possible in the sparse regime. Both scopes are retained. If \(M\subseteq V\) is Maker's claimed set, a possible majority objective is \(\exists A\in\mathcal E\) with \(|M\cap A|>|A|/2\); a possible biased-surplus expression is
\[
 |M\cap A|-\frac{a}{a+b}|A|\geq s(A).
\]
These are examples of data that a precise objective must specify, not objectives selected by the source. Bias may change the appropriate baseline, and strict versus weak inequalities and a required surplus amount matter.

Sparse is not defined by an edge count, uniformity, vertex degree, codegree or another parameter, and the desired positive statement could depend on those choices and the chosen objective. One must specify the hypergraph class, asymptotic regime and winning criterion before a concrete strategy theorem can be stated. No assumption that a complete-occupation theorem automatically transfers to a majority game is made.

The literature triage checked 2026-08-21 reports the general biased Meta-Conjecture and useful sparse-case results broadly unresolved, with later surveys still treating the biased case as open. This dated description, the open label and reported source verification remain supplied metadata, not independently checked current conclusions.
\subsection{Short English statement}
Extend the source's missing weak-win Meta-Conjecture to biased hypergraph games, and find meaningful guarantees in a specified sparse regime, possibly for majority or surplus goals. The game rules, conjecture and quantitative sparse/objective conditions must be supplied before a definite theorem can be chosen.
\subsection{Sources}
\begin{itemize}
\item Local UnsolvedMath registry, permanent ID 30000731, source code catalogue entry, ``Biased Weak-Win Games on Sparse Hypergraphs,'' including the original two questions, Beck attribution in the cleaned statement and its majority/surplus refinement.
\item \emph{Mini-Workshop: Positional Games (2008)}, Oberwolfach Reports, pp.\ 1073--1114, DOI 10.4171/owr/2007/20: \url{https://doi.org/10.4171/owr/2007/20}. The citation uses 2008, while the DOI and proposed-year metadata use 2007.
\item Supplied reference ``biased Meta-Conjecture survey'': \url{https://doi.org/10.1017/CBO9780511735202.011}.
\item Literature triage checked 2026-08-21; record updated 2026-09-16. No external source has been visited or independently checked here.
\end{itemize}
\end{document}
